\bibitem{support-type-one-thorne}
J. A. Thorne, \emph{On the vanishing of adjoint Bloch--Kato Selmer
groups of irreducible automorphic Galois representations},
Pure and Applied Mathematics Quarterly \textbf{18} (2022), no.~5,
Theorem~6.1.  DOI: \texttt{10.4310/PAMQ.2022.v18.n5.a5}.
The accepted manuscript is \texttt{arXiv:2207.04925v1}; see
Theorem~A on p.~3, the adjoint notation on p.~4, and
Theorem~6.1 and the multiplier in \S6.3 on pp.~24--25.
\url{https://arxiv.org/abs/2207.04925v1}.

\bibitem{support-type-one-source}
L. Alp\"oge, pinned source \texttt{jsmorph/rh}, revision
\texttt{ecd09ff9b4aaf6c08f2b5ef17dc49ca9648a77de}, under
\path{Development/AutomorphicArithmetic/AutomorphicVarieties}.
The separate decomposition-group and inertia elements occur explicitly
in \path{UnitaryThreeDeformations/IrrOpen.lean}, 275--279;
its residual trace argument is at 258--269 and 315--338.
The original local points and rigidity are in
\path{UnitaryThreeGroups/GRepresents.lean}, 128--141.
The normalized arithmetic-point automorphy contract is in
\path{CurveDVRSpecialization/RankTwoAutomorphyDescendsFromAbelianCMToRat.lean},
16--20; its unchanged Hecke-classicality input is in
\path{WeilGroupCharacters/ArithmeticPointClassicality.lean}, 99--101.
The exact adjoint coefficient module and the local tangent consumer
are in \path{PadicHodgeTheory/DualCr.lean}, 36--62, and
\path{UnitaryThreeGroups/TanCocIsBK.lean}, 187--222.
Authenticated complete-file hashes and bounded inspected ranges are
recorded in
\path{certificates/selmer-support-continuation-type-one-rigidity.json}.
